\section{Related Work}

\paragraph{LLM judges and human alignment.}
LLM judges are now widely used to score open-ended outputs, with human agreement as the dominant validation criterion \citep{Liu2023GEval,Zheng2023MTBench,Kocmi2024HumanVsLLM}.
Large-scale studies show this agreement can vary substantially across tasks and prompts, motivating stronger validity checks before replacing human raters \citep{Kocmi2024HumanVsLLM,Li2024TrustLLMJudgments}.
Our work keeps human alignment as an anchor but introduces property-level constraints that can be audited and re-run cheaply.

\paragraph{Reliability and robustness of automated evaluation.}
Prior work improves judge reliability via structured prompting and multi-agent judging pipelines \citep{Liu2023GEval,Chan2024ChatEval}.
These methods reduce some noise sources, but they do not by themselves guarantee core evaluator behaviors such as identity or monotone response to controlled degradation.
We complement these approaches with explicit falsification tests on perturbation ladders.

\paragraph{Code and formal-verification evaluation.}
In code tasks, judge-based metrics are useful when exact execution-based grading is unavailable or incomplete \citep{TongZhang2024CodeJudge,Zhuo2024ICEScore}.
In formal verification, many benchmarks provide fixed target theorems, allowing proof success to dominate evaluation.
Agentic verification differs because systems must also infer \emph{what should be proved}.
This makes semantic theorem/specification quality a central measurement problem and motivates certification of the judge itself.

\paragraph{Scalable oversight and AI measurement science.}
Scalable oversight work argues that automated evaluators must be calibrated against human standards and stress-tested under realistic failure modes \citep{bowman2022scalable,zhou2025graders}.
Our protocol contributes one concrete mechanism for that agenda: a reusable validity layer for judge selection that can be recalibrated across model updates and distribution shifts.
